\documentclass[aspectratio=169,8pt]{beamer}
\usetheme{Madrid}
\usepackage[utf8]{inputenc}
\usepackage[T1]{fontenc}
\usepackage{amsmath,amssymb}
\usepackage{booktabs}
\usepackage{enumitem}
\setlist[enumerate]{label=\Alph*.,leftmargin=*,itemsep=1pt,topsep=2pt}
\setlist[itemize]{leftmargin=*,itemsep=1pt,topsep=2pt}
\setbeamerfont{block body}{size=\tiny}
\setbeamerfont{block title}{size=\scriptsize}
\setbeamerfont{frametitle}{size=\footnotesize}
\setbeamerfont{normal text}{size=\tiny}
\setbeamertemplate{navigation symbols}{}
\setbeamertemplate{itemize item}{\tiny$\bullet$}
\setbeamertemplate{itemize subitem}{\tiny$\circ$}

\title[GARP RAI]{GARP Risk \& AI --- Q551--Q650}
\subtitle{Unofficial Exam Prep}
\author{Unofficial}
\date{\today}

\begin{document}
	
	\begin{frame}
		\titlepage
	\end{frame}
	
	% =================================================================
	\section{Advanced AI Risk Taxonomy}
	% =================================================================
	
	\begin{frame}{Q551 --- Model Risk vs Technology Risk}
		\begin{block}{Question}
			How does model risk differ from technology risk in an AI governance framework?
			\begin{enumerate}
				\item They are identical.
				\item Model risk focuses on potential adverse consequences from model errors, misuse, or limitations; technology risk focuses on IT infrastructure, systems, and operational failures.
				\item Model risk applies only to AI; technology risk applies only to statistics.
				\item Neither is regulated.
			\end{enumerate}
		\end{block}
		\begin{exampleblock}{Answer: B}\end{exampleblock}
		\begin{block}{Explanation}
			Model risk relates to model-based decisions; technology risk covers systems and infrastructure. Options A, C, and D misstate the distinction.
		\end{block}
	\end{frame}
	
	\begin{frame}{Q552 --- AI Risk Register}
		\begin{block}{Question}
			Which of the following is the primary purpose of an AI risk register?
			\begin{enumerate}
				\item List every AI model in the firm.
				\item Track identified AI risks, their owners, mitigation actions, and remediation status.
				\item Replace the model inventory.
				\item Document marketing strategy.
			\end{enumerate}
		\end{block}
		\begin{exampleblock}{Answer: B}\end{exampleblock}
		\begin{block}{Explanation}
			The AI risk register tracks identified risks, owners, mitigations, and status. Options A, C, and D are different artifacts or unrelated.
		\end{block}
	\end{frame}
	
	\begin{frame}{Q553 --- AI Risk Appetite}
		\begin{block}{Question}
			Why does an organisation need an AI risk appetite?
			\begin{enumerate}
				\item To replace regulation.
				\item To define the level and types of AI risk the organisation is willing to accept, guiding investment and control decisions.
				\item To eliminate all AI risk.
				\item To market AI products.
			\end{enumerate}
		\end{block}
		\begin{exampleblock}{Answer: B}\end{exampleblock}
		\begin{block}{Explanation}
			AI risk appetite defines acceptable risk to guide decisions. Options A, C, and D misstate its role.
		\end{block}
	\end{frame}
	
	\begin{frame}{Q554 --- Enterprise-Wide AI Governance}
		\begin{block}{Question}
			What is the key purpose of enterprise-wide AI governance?
			\begin{enumerate}
				\item Local optimisation of one team.
				\item Consistent policies, oversight, and controls applied across all AI initiatives, aligned to overall business strategy and risk appetite.
				\item Eliminating AI use.
				\item Outsourcing AI to vendors.
			\end{enumerate}
		\end{block}
		\begin{exampleblock}{Answer: B}\end{exampleblock}
		\begin{block}{Explanation}
			Enterprise AI governance provides consistent policies and oversight. Options A, C, and D are inconsistent with the concept.
		\end{block}
	\end{frame}
	
	\begin{frame}{Q555 --- AI Risk Tolerance Triggers}
		\begin{block}{Question}
			Which event should trigger an escalation in an AI governance framework?
			\begin{enumerate}
				\item Routine model retraining.
				\item Material breach of AI risk tolerance, such as a serious fairness violation or a regulatory finding.
				\item Minor UI changes.
				\item Documentation updates.
			\end{enumerate}
		\end{block}
		\begin{exampleblock}{Answer: B}\end{exampleblock}
		\begin{block}{Explanation}
			Material breaches require escalation. Options A, C, and D are routine activities.
		\end{block}
	\end{frame}
	
	\begin{frame}{Q556 --- AI Risk Reporting Frequency}
		\begin{block}{Question}
			How often should AI risk reporting occur to senior management?
			\begin{enumerate}
				\item Only annually.
				\item Regularly (e.g., quarterly) and ad hoc for significant events.
				\item Never.
				\item Only at year-end.
			\end{enumerate}
		\end{block}
		\begin{exampleblock}{Answer: B}\end{exampleblock}
		\begin{block}{Explanation}
			Reporting is regular and ad hoc. Options A, C, and D are insufficient.
		\end{block}
	\end{frame}
	
	\begin{frame}{Q557 --- Model Inventory vs AI Portfolio}
		\begin{block}{Question}
			How does an AI portfolio differ from a model inventory?
			\begin{enumerate}
				\item They are identical.
				\item A portfolio is a strategic view of AI capabilities and initiatives; the inventory is an operational registry of deployed models.
				\item Portfolio is only for vendors.
				\item Inventory is only for retired models.
			\end{enumerate}
		\end{block}
		\begin{exampleblock}{Answer: B}\end{exampleblock}
		\begin{block}{Explanation}
			Portfolio is strategic; inventory is operational. Options A, C, and D misstate.
		\end{block}
	\end{frame}
	
	\begin{frame}{Q558 --- AI Governance Board Reporting}
		\begin{block}{Question}
			Which is NOT typically included in board-level AI reporting?
			\begin{enumerate}
				\item Aggregate AI risk exposure.
				\item Fairness and bias metrics.
				\item Detailed source code of every model.
				\item Material incidents and remediation plans.
			\end{enumerate}
		\end{block}
		\begin{exampleblock}{Answer: C}\end{exampleblock}
		\begin{block}{Explanation}
			Board reports are strategic and aggregate; source code is not appropriate. Options A, B, and D are typically included.
		\end{block}
	\end{frame}
	
	\begin{frame}{Q559 --- Systemic AI Risk}
		\begin{block}{Question}
			What is systemic AI risk?
			\begin{enumerate}
				\item A minor operational issue.
				\item Risk that a widespread AI failure or correlated AI behaviour could affect multiple firms or the broader financial system.
				\item Only internal risk.
				\item Irrelevant to regulation.
			\end{enumerate}
		\end{block}
		\begin{exampleblock}{Answer: B}\end{exampleblock}
		\begin{block}{Explanation}
			Systemic AI risk extends beyond single-firm boundaries. Options A, C, and D understate or dismiss.
		\end{block}
	\end{frame}
	
	\begin{frame}{Q560 --- AI Risk Taxonomy Categories}
		\begin{block}{Question}
			Which is a common category in AI risk taxonomies?
			\begin{enumerate}
				\item Technical risk.
				\item Ethical risk.
				\item Operational risk.
				\item All of the above.
			\end{enumerate}
		\end{block}
		\begin{exampleblock}{Answer: D}\end{exampleblock}
		\begin{block}{Explanation}
			AI risk taxonomies span technical, ethical, operational, and other categories.
		\end{block}
	\end{frame}
	
	% =================================================================
	\section{AI Model Risk in Financial Services}
	% =================================================================
	
	\begin{frame}{Q561 --- Credit Scoring Model Risk}
		\begin{block}{Question}
			Which is the most relevant model risk for an AI-based credit scoring system?
			\begin{enumerate}
				\item Colour of the user interface.
				\item Potential unfairness, inaccuracy, and adverse regulatory consequences from biased or poorly performing predictions.
				\item Storage cost only.
				\item Marketing effectiveness only.
			\end{enumerate}
		\end{block}
		\begin{exampleblock}{Answer: B}\end{exampleblock}
		\begin{block}{Explanation}
			Credit scoring model risk includes fairness, accuracy, and regulatory risk. Options A, C, and D are secondary or irrelevant.
		\end{block}
	\end{frame}
	
	\begin{frame}{Q562 --- Fraud Detection Model Risk}
		\begin{block}{Question}
			For an AI fraud detection model, which is the most critical model risk?
			\begin{enumerate}
				\item Overfitting to historical fraud patterns, missing novel fraud schemes.
				\item Storage cost.
				\item User interface design.
				\item Training time.
			\end{enumerate}
		\end{block}
		\begin{exampleblock}{Answer: A}\end{exampleblock}
		\begin{block}{Explanation}
			Overfitting to past patterns can miss novel fraud. Options B, C, and D are less critical.
		\end{block}
	\end{frame}
	
	\begin{frame}{Q563 --- AML Model Risk}
		\begin{block}{Question}
			Which is a key risk for an AI anti-money-laundering (AML) model?
			\begin{enumerate}
				\item Low alert volume.
				\item High false-positive rate overwhelming compliance teams and potential missed true positives, both creating regulatory and operational risk.
				\item Storage cost.
				\item Marketing.
			\end{enumerate}
		\end{block}
		\begin{exampleblock}{Answer: B}\end{exampleblock}
		\begin{block}{Explanation}
			AML risk includes excessive false positives and missed true positives. Options A, C, and D are secondary.
		\end{block}
	\end{frame}
	
	\begin{frame}{Q564 --- Market Risk Model}
		\begin{block}{Question}
			Which is a typical risk for an AI-enhanced market risk model?
			\begin{enumerate}
				\item Training on non-stationary data leading to unstable predictions.
				\item Data storage.
				\item Logo design.
				\item Employee satisfaction.
			\end{enumerate}
		\end{block}
		\begin{exampleblock}{Answer: A}\end{exampleblock}
		\begin{block}{Explanation}
			Non-stationarity of financial markets is a key challenge. Options B, C, and D are unrelated.
		\end{block}
	\end{frame}
	
	\begin{frame}{Q565 --- Insurance Underwriting AI}
		\begin{block}{Question}
			Which is a specific concern for AI underwriting in insurance?
			\begin{enumerate}
				\item Ethical and legal issues from using alternative data sources that may be considered sensitive or discriminatory.
				\item Faster approvals.
				\item Lower costs.
				\item Better customer experience.
			\end{enumerate}
		\end{block}
		\begin{exampleblock}{Answer: A}\end{exampleblock}
		\begin{block}{Explanation}
			Alternative data raises ethical and legal concerns. Options B, C, and D are benefits, not concerns.
		\end{block}
	\end{frame}
	
	\begin{frame}{Q566 --- Model Fairness in Lending}
		\begin{block}{Question}
			Which is a recognised fairness metric used in lending AI?
			\begin{enumerate}
				\item Demographic parity.
				\item Storage cost parity.
				\item Latency parity.
				\item Colour parity.
			\end{enumerate}
		\end{block}
		\begin{exampleblock}{Answer: A}\end{exampleblock}
		\begin{block}{Explanation}
			Demographic parity is a standard fairness metric. Options B, C, and D are not.
		\end{block}
	\end{frame}
	
	\begin{frame}{Q567 --- Algorithmic Trading Model Risk}
		\begin{block}{Question}
			Which is a key risk of algorithmic trading models?
			\begin{enumerate}
				\item Flash-crash amplification, unintended feedback loops, and market impact from correlated strategies.
				\item Storage capacity.
				\item Logo design.
				\item Employee benefits.
			\end{enumerate}
		\end{block}
		\begin{exampleblock}{Answer: A}\end{exampleblock}
		\begin{block}{Explanation}
			Algorithmic trading risks include flash crashes and correlated behaviour. Options B, C, and D are unrelated.
		\end{block}
	\end{frame}
	
	\begin{frame}{Q568 --- Chatbot Risk in Banking}
		\begin{block}{Question}
			Which is the most significant risk from deploying a GenAI chatbot in banking?
			\begin{enumerate}
				\item Generating incorrect or non-compliant advice to customers, creating legal and reputational risk.
				\item Faster response times.
				\item Cost savings.
				\item Customer satisfaction.
			\end{enumerate}
		\end{block}
		\begin{exampleblock}{Answer: A}\end{exampleblock}
		\begin{block}{Explanation}
			Advice quality and compliance are central concerns. Options B, C, and D are benefits.
		\end{block}
	\end{frame}
	
	\begin{frame}{Q569 --- Model Risk from Alternative Data}
		\begin{block}{Question}
			Which is a key governance concern when using alternative data in financial AI?
			\begin{enumerate}
				\item Legal right to use the data, privacy, and potential for discriminatory impact.
				\item Storage cost.
				\item Model speed.
				\item Logo design.
			\end{enumerate}
		\end{block}
		\begin{exampleblock}{Answer: A}\end{exampleblock}
		\begin{block}{Explanation}
			Legal rights, privacy, and fairness are key governance concerns. Options B, C, and D are secondary.
		\end{block}
	\end{frame}
	
	\begin{frame}{Q570 --- Model Risk in AML/KYC}
		\begin{block}{Question}
			What is a specific AI risk in AML/KYC?
			\begin{enumerate}
				\item Customer friction from over-flagging and regulatory exposure from under-flagging, requiring careful threshold calibration.
				\item Storage cost.
				\item Logo design.
				\item Training time.
			\end{enumerate}
		\end{block}
		\begin{exampleblock}{Answer: A}\end{exampleblock}
		\begin{block}{Explanation}
			Threshold calibration affects both customer experience and regulatory compliance. Options B, C, and D are less critical.
		\end{block}
	\end{frame}
	
	% =================================================================
	\section{AI Model Lifecycle and Change Management}
	% =================================================================
	
	\begin{frame}{Q571 --- Model Lifecycle Stages}
		\begin{block}{Question}
			Which is the correct order of model lifecycle stages?
			\begin{enumerate}
				\item Development, validation, deployment, monitoring, retirement.
				\item Deployment, development, monitoring, validation, retirement.
				\item Monitoring, development, retirement, validation, deployment.
				\item Retirement, development, deployment, monitoring, validation.
			\end{enumerate}
		\end{block}
		\begin{exampleblock}{Answer: A}\end{exampleblock}
		\begin{block}{Explanation}
			Development $\to$ validation $\to$ deployment $\to$ monitoring $\to$ retirement. Options B, C, and D reorder incorrectly.
		\end{block}
	\end{frame}
	
	\begin{frame}{Q572 --- Validation Before Deployment}
		\begin{block}{Question}
			Why must validation occur before deployment?
			\begin{enumerate}
				\item Regulatory convenience.
				\item To independently confirm the model meets requirements, performs as expected, and is fit for purpose before being used in production.
				\item To reduce cost.
				\item To increase speed.
			\end{enumerate}
		\end{block}
		\begin{exampleblock}{Answer: B}\end{exampleblock}
		\begin{block}{Explanation}
			Pre-deployment validation confirms fitness. Options A, C, and D are secondary.
		\end{block}
	\end{frame}
	
	\begin{frame}{Q573 --- Post-Deployment Monitoring}
		\begin{block}{Question}
			Which activity is part of post-deployment monitoring?
			\begin{enumerate}
				\item Continuous tracking of model performance, data drift, and input integrity, with alerts for threshold breaches.
				\item Only training.
				\item Only validation.
				\item Only code review.
			\end{enumerate}
		\end{block}
		\begin{exampleblock}{Answer: A}\end{exampleblock}
		\begin{block}{Explanation}
			Post-deployment monitoring tracks performance and drift. Options B, C, and D are pre-deployment.
		\end{block}
	\end{frame}
	
	\begin{frame}{Q574 --- Recalibration vs Redevelopment}
		\begin{block}{Question}
			Recalibration is preferred to redevelopment when:
			\begin{enumerate}
				\item The model structure is wrong.
				\item The model structure remains valid but parameters need updating using new data.
				\item The model is retired.
				\item The model has never been validated.
			\end{enumerate}
		\end{block}
		\begin{exampleblock}{Answer: B}\end{exampleblock}
		\begin{block}{Explanation}
			Recalibration updates parameters; structural issues require redevelopment. Options A, C, and D are different situations.
		\end{block}
	\end{frame}
	
	\begin{frame}{Q575 --- Model Deprecation}
		\begin{block}{Question}
			What is model deprecation?
			\begin{enumerate}
				\item Immediate retirement.
				\item The formal marking of a model as no longer recommended for use, typically while a replacement is being prepared.
				\item Model creation.
				\item Model validation.
			\end{enumerate}
		\end{block}
		\begin{exampleblock}{Answer: B}\end{exampleblock}
		\begin{block}{Explanation}
			Deprecation flags models for phase-out. Options A, C, and D describe different actions.
		\end{block}
	\end{frame}
	
	\begin{frame}{Q576 --- Model Change Impact Assessment}
		\begin{block}{Question}
			What does a model change impact assessment evaluate?
			\begin{enumerate}
				\item Only cost.
				\item How a proposed change affects model outputs, downstream users, risk profile, and regulatory compliance.
				\item Only speed.
				\item Only storage.
			\end{enumerate}
		\end{block}
		\begin{exampleblock}{Answer: B}\end{exampleblock}
		\begin{block}{Explanation}
			Impact assessment covers outputs, downstream, risk, and compliance. Options A, C, and D are too narrow.
		\end{block}
	\end{frame}
	
	\begin{frame}{Q577 --- Model Versioning}
		\begin{block}{Question}
			Why is model versioning important?
			\begin{enumerate}
				\item To reduce storage.
				\item To track distinct model iterations, enabling reproducibility, audit, and controlled rollback if needed.
				\item To reduce cost.
				\item To speed up training.
			\end{enumerate}
		\end{block}
		\begin{exampleblock}{Answer: B}\end{exampleblock}
		\begin{block}{Explanation}
			Versioning supports reproducibility, audit, and rollback. Options A, C, and D are secondary.
		\end{block}
	\end{frame}
	
	\begin{frame}{Q578 --- Champion/Challenger Framework}
		\begin{block}{Question}
			What is the champion/challenger framework?
			\begin{enumerate}
				\item A marketing strategy.
				\item A framework where the current production model (champion) is compared against alternative models (challengers) to evaluate whether a switch would improve performance.
				\item A training technique.
				\item A validation report.
			\end{enumerate}
		\end{block}
		\begin{exampleblock}{Answer: B}\end{exampleblock}
		\begin{block}{Explanation}
			Champion/challenger compares production models to alternatives. Options A, C, and D are unrelated.
		\end{block}
	\end{frame}
	
	\begin{frame}{Q579 --- Shadow Deployment}
		\begin{block}{Question}
			What is shadow deployment of a model?
			\begin{enumerate}
				\item Deploying without monitoring.
				\item Running a new model alongside production without using its outputs, to evaluate performance in real conditions.
				\item Deploying to production only.
				\item Retiring a model.
			\end{enumerate}
		\end{block}
		\begin{exampleblock}{Answer: B}\end{exampleblock}
		\begin{block}{Explanation}
			Shadow deployment evaluates models in parallel without affecting production. Options A, C, and D misstate.
		\end{block}
	\end{frame}
	
	\begin{frame}{Q580 --- A/B Testing of Models}
		\begin{block}{Question}
			What is A/B testing in the AI deployment context?
			\begin{enumerate}
				\item Randomly assigning users to two model versions and comparing outcomes.
				\item Retiring both models.
				\item Doubling training data.
				\item Storing data twice.
			\end{enumerate}
		\end{block}
		\begin{exampleblock}{Answer: A}\end{exampleblock}
		\begin{block}{Explanation}
			A/B testing compares model versions via controlled experiments. Options B, C, and D are unrelated.
		\end{block}
	\end{frame}
	
	% =================================================================
	\section{AI Incident Management}
	% =================================================================
	
	\begin{frame}{Q581 --- AI Incident Response}
		\begin{block}{Question}
			Which is a key element of an AI incident response plan?
			\begin{enumerate}
				\item Ignoring incidents.
				\item Defined escalation paths, roles, playbooks, and communication plans for handling AI-related incidents.
				\item Only post-mortems.
				\item Only IT involvement.
			\end{enumerate}
		\end{block}
		\begin{exampleblock}{Answer: B}\end{exampleblock}
		\begin{block}{Explanation}
			An incident response plan includes escalation, roles, playbooks, and communication. Options A, C, and D are insufficient.
		\end{block}
	\end{frame}
	
	\begin{frame}{Q582 --- AI Incident Root Cause}
		\begin{block}{Question}
			An AI model begins producing biased outputs. Which is the most likely root cause?
			\begin{enumerate}
				\item Data drift or newly biased data.
				\item Server failure.
				\item Power outage.
				\item UI redesign.
			\end{enumerate}
		\end{block}
		\begin{exampleblock}{Answer: A}\end{exampleblock}
		\begin{block}{Explanation}
			Bias in outputs typically stems from data or model development. Options B, C, and D are unrelated.
		\end{block}
	\end{frame}
	
	\begin{frame}{Q583 --- Post-Incident Review}
		\begin{block}{Question}
			What is the purpose of a post-incident review for AI incidents?
			\begin{enumerate}
				\item Assign blame.
				\item Learn from the incident, improve controls and processes, and prevent recurrence.
				\item End investigation.
				\item Skip documentation.
			\end{enumerate}
		\end{block}
		\begin{exampleblock}{Answer: B}\end{exampleblock}
		\begin{block}{Explanation}
			Post-incident reviews focus on learning and prevention. Options A, C, and D are contrary.
		\end{block}
	\end{frame}
	
	\begin{frame}{Q584 --- Incident Communication}
		\begin{block}{Question}
			Why is communication important during an AI incident?
			\begin{enumerate}
				\item It is not.
				\item Clear communication with regulators, customers, and stakeholders maintains trust and meets regulatory obligations.
				\item Only internal communication matters.
				\item Only external communication matters.
			\end{enumerate}
		\end{block}
		\begin{exampleblock}{Answer: B}\end{exampleblock}
		\begin{block}{Explanation}
			Communication with all stakeholders is essential. Options A, C, and D are incomplete or wrong.
		\end{block}
	\end{frame}
	
	\begin{frame}{Q585 --- AI Incident Classification}
		\begin{block}{Question}
			How are AI incidents typically classified?
			\begin{enumerate}
				\item By colour.
				\item By severity and impact, considering factors such as harm to customers, financial loss, and regulatory exposure.
				\item Randomly.
				\item By alphabet.
			\end{enumerate}
		\end{block}
		\begin{exampleblock}{Answer: B}\end{exampleblock}
		\begin{block}{Explanation}
			Incidents are classified by severity and impact. Options A, C, and D are not standard.
		\end{block}
	\end{frame}
	
	\begin{frame}{Q586 --- Regulatory Notification}
		\begin{block}{Question}
			When a serious AI incident affects customers, what is typically required?
			\begin{enumerate}
				\item No action.
				\item Prompt notification to regulators where required and remediation plans.
				\item Only internal review.
				\item Silence.
			\end{enumerate}
		\end{block}
		\begin{exampleblock}{Answer: B}\end{exampleblock}
		\begin{block}{Explanation}
			Serious incidents may trigger regulatory notification. Options A, C, and D violate governance.
		\end{block}
	\end{frame}
	
	\begin{frame}{Q587 --- Incident Metrics}
		\begin{block}{Question}
			Which metric helps track AI incident management effectiveness?
			\begin{enumerate}
				\item Total storage used.
				\item Time to detection, time to remediation, recurrence rate, and impact.
				\item Marketing reach.
				\item UI clicks.
			\end{enumerate}
		\end{block}
		\begin{exampleblock}{Answer: B}\end{exampleblock}
		\begin{block}{Explanation}
			Detection, remediation, recurrence, and impact are key metrics. Options A, C, and D are unrelated.
		\end{block}
	\end{frame}
	
	\begin{frame}{Q588 --- Red Teaming}
		\begin{block}{Question}
			What is red teaming in AI?
			\begin{enumerate}
				\item A marketing exercise.
				\item Adversarial testing by a dedicated team to uncover vulnerabilities and unintended behaviours before deployment.
				\item A financial audit.
				\item A user survey.
			\end{enumerate}
		\end{block}
		\begin{exampleblock}{Answer: B}\end{exampleblock}
		\begin{block}{Explanation}
			Red teaming proactively tests for vulnerabilities. Options A, C, and D misstate.
		\end{block}
	\end{frame}
	
	\begin{frame}{Q589 --- AI Incident Data}
		\begin{block}{Question}
			Why is logging data important for AI incident management?
			\begin{enumerate}
				\item To increase storage.
				\item To reconstruct events, identify root causes, and support regulatory reporting.
				\item To slow down systems.
				\item To increase costs.
			\end{enumerate}
		\end{block}
		\begin{exampleblock}{Answer: B}\end{exampleblock}
		\begin{block}{Explanation}
			Logging supports root-cause analysis and reporting. Options A, C, and D are incorrect or secondary.
		\end{block}
	\end{frame}
	
	\begin{frame}{Q590 --- Continuous Improvement}
		\begin{block}{Question}
			What is the purpose of continuous improvement in AI governance?
			\begin{enumerate}
				\item To increase costs.
				\item To iterate on governance frameworks and controls based on incidents, audits, and evolving best practices.
				\item To avoid change.
				\item To slow down deployment.
			\end{enumerate}
		\end{block}
		\begin{exampleblock}{Answer: B}\end{exampleblock}
		\begin{block}{Explanation}
			Continuous improvement ensures governance keeps pace with change. Options A, C, and D are contrary.
		\end{block}
	\end{frame}
	
	% =================================================================
	\section{AI Governance Maturity}
	% =================================================================
	
	\begin{frame}{Q591 --- AI Governance Maturity Model}
		\begin{block}{Question}
			A typical AI governance maturity model includes levels such as:
			\begin{enumerate}
				\item Ad hoc, developing, defined, managed, optimised.
				\item Red, blue, green.
				\item Only one level.
				\item No levels.
			\end{enumerate}
		\end{block}
		\begin{exampleblock}{Answer: A}\end{exampleblock}
		\begin{block}{Explanation}
			Maturity models typically progress from ad hoc to optimised. Options B, C, and D are not standard.
		\end{block}
	\end{frame}
	
	\begin{frame}{Q592 --- AI Governance Gaps}
		\begin{block}{Question}
			Which is a common gap in AI governance frameworks?
			\begin{enumerate}
				\item Extensive controls.
				\item Missing governance for GenAI, alternative data, or third-party models.
				\item Complete coverage.
				\item Perfect alignment.
			\end{enumerate}
		\end{block}
		\begin{exampleblock}{Answer: B}\end{exampleblock}
		\begin{block}{Explanation}
			Common gaps include GenAI, alternative data, and vendor models. Options A, C, and D are not gaps.
		\end{block}
	\end{frame}
	
	\begin{frame}{Q593 --- Governance Roadmap}
		\begin{block}{Question}
			What is the purpose of an AI governance roadmap?
			\begin{enumerate}
				\item A marketing document.
				\item A phased plan outlining how the organisation will build, enhance, and sustain AI governance capabilities over time.
				\item A single policy.
				\item A code repository.
			\end{enumerate}
		\end{block}
		\begin{exampleblock}{Answer: B}\end{exampleblock}
		\begin{block}{Explanation}
			A roadmap is a phased plan. Options A, C, and D misstate.
		\end{block}
	\end{frame}
	
	\begin{frame}{Q594 --- AI Governance KPIs}
		\begin{block}{Question}
			Which is an appropriate KPI for AI governance?
			\begin{enumerate}
				\item Number of employees.
				\item Percentage of models validated on schedule.
				\item Office size.
				\item Marketing spend.
			\end{enumerate}
		\end{block}
		\begin{exampleblock}{Answer: B}\end{exampleblock}
		\begin{block}{Explanation}
			Validation timeliness is a governance KPI. Options A, C, and D are unrelated.
		\end{block}
	\end{frame}
	
	\begin{frame}{Q595 --- AI Governance Training}
		\begin{block}{Question}
			Which audience requires AI governance training?
			\begin{enumerate}
				\item Only data scientists.
				\item A broad set including executives, developers, validators, business users, compliance, and internal audit.
				\item Only IT.
				\item Only external auditors.
			\end{enumerate}
		\end{block}
		\begin{exampleblock}{Answer: B}\end{exampleblock}
		\begin{block}{Explanation}
			AI governance training spans many roles. Options A, C, and D are too narrow.
		\end{block}
	\end{frame}
	
	\begin{frame}{Q596 --- Governance Data Collection}
		\begin{block}{Question}
			Why is data collection for AI governance metrics important?
			\begin{enumerate}
				\item To meet marketing goals.
				\item To enable measurement, benchmarking, and continuous improvement of AI governance practices.
				\item To increase costs.
				\item To slow operations.
			\end{enumerate}
		\end{block}
		\begin{exampleblock}{Answer: B}\end{exampleblock}
		\begin{block}{Explanation}
			Data supports measurement and improvement. Options A, C, and D are incorrect.
		\end{block}
	\end{frame}
	
	\begin{frame}{Q597 --- AI Governance Certification}
		\begin{block}{Question}
			What is the purpose of AI governance certifications (e.g., ISO/IEC 42001)?
			\begin{enumerate}
				\item Marketing only.
				\item To demonstrate externally that the organisation meets recognised standards for AI governance.
				\item To replace regulations.
				\item To lower costs.
			\end{enumerate}
		\end{block}
		\begin{exampleblock}{Answer: B}\end{exampleblock}
		\begin{block}{Explanation}
			Certifications demonstrate conformance with standards. Options A, C, and D are not the primary purpose.
		\end{block}
	\end{frame}
	
	\begin{frame}{Q598 --- Board AI Literacy}
		\begin{block}{Question}
			Why is board AI literacy important?
			\begin{enumerate}
				\item To increase headcount.
				\item To enable effective oversight of AI strategy, risk, and governance at the highest level.
				\item To reduce costs.
				\item To speed deployment.
			\end{enumerate}
		\end{block}
		\begin{exampleblock}{Answer: B}\end{exampleblock}
		\begin{block}{Explanation}
			Board AI literacy enables effective oversight. Options A, C, and D are not the primary reason.
		\end{block}
	\end{frame}
	
	\begin{frame}{Q599 --- AI Ethics Committee Integration}
		\begin{block}{Question}
			How should an AI ethics committee integrate with the model risk function?
			\begin{enumerate}
				\item No integration is needed.
				\item Through coordinated review of high-risk AI use cases, sharing of insights on fairness and ethics, and joint escalation to senior management.
				\item Only through annual meetings.
				\item Only through external advisors.
			\end{enumerate}
		\end{block}
		\begin{exampleblock}{Answer: B}\end{exampleblock}
		\begin{block}{Explanation}
			Coordination between ethics and model risk is essential. Options A, C, and D are insufficient.
		\end{block}
	\end{frame}
	
	\begin{frame}{Q600 --- Governance Evolution}
		\begin{block}{Question}
			Why must AI governance frameworks evolve?
			\begin{enumerate}
				\item To increase cost.
				\item Because AI technologies, use cases, and regulations change rapidly, requiring ongoing updates to remain effective.
				\item To slow operations.
				\item To reduce transparency.
			\end{enumerate}
		\end{block}
		\begin{exampleblock}{Answer: B}\end{exampleblock}
		\begin{block}{Explanation}
			Rapid change in AI and regulation requires evolving governance. Options A, C, and D contradict this.
		\end{block}
	\end{frame}
	
	% =================================================================
	\section{Ethics in GenAI}
	% =================================================================
	
	\begin{frame}{Q601 --- GenAI and Autonomy}
		\begin{block}{Question}
			How can GenAI undermine human autonomy?
			\begin{enumerate}
				\item By providing choices.
				\item By generating persuasive, tailored content that manipulates users without their awareness.
				\item By explaining decisions.
				\item By slowing systems.
			\end{enumerate}
		\end{block}
		\begin{exampleblock}{Answer: B}\end{exampleblock}
		\begin{block}{Explanation}
			Manipulative content undermines autonomy. Options A, C, and D respect or are unrelated to autonomy.
		\end{block}
	\end{frame}
	
	\begin{frame}{Q602 --- GenAI and Bias}
		\begin{block}{Question}
			How does GenAI amplify bias?
			\begin{enumerate}
				\item By generating content that reflects biases present in training data, at scale.
				\item By reducing accuracy.
				\item By slowing training.
				\item By increasing costs.
			\end{enumerate}
		\end{block}
		\begin{exampleblock}{Answer: A}\end{exampleblock}
		\begin{block}{Explanation}
			GenAI can amplify training-data bias at scale. Options B, C, and D are unrelated.
		\end{block}
	\end{frame}
	
	\begin{frame}{Q603 --- GenAI and Truth}
		\begin{block}{Question}
			Why is GenAI a challenge for truth and accountability?
			\begin{enumerate}
				\item It only outputs truth.
				\item It generates plausible but sometimes incorrect content, making accountability for misinformation difficult.
				\item It has no training data.
				\item It never hallucinates.
			\end{enumerate}
		\end{block}
		\begin{exampleblock}{Answer: B}\end{exampleblock}
		\begin{block}{Explanation}
			Hallucinations and plausible output complicate accountability. Options A, C, and D are incorrect.
		\end{block}
	\end{frame}
	
	\begin{frame}{Q604 --- GenAI and Consent}
		\begin{block}{Question}
			What consent issues arise with GenAI training data?
			\begin{enumerate}
				\item None.
				\item Data subjects whose content was used in training may not have consented, raising ethical and legal concerns.
				\item Consent is automatically granted.
				\item Consent is not required.
			\end{enumerate}
		\end{block}
		\begin{exampleblock}{Answer: B}\end{exampleblock}
		\begin{block}{Explanation}
			Training data consent is a critical concern. Options A, C, and D are incorrect.
		\end{block}
	\end{frame}
	
	\begin{frame}{Q605 --- GenAI and Attribution}
		\begin{block}{Question}
			Why is attribution important in GenAI?
			\begin{enumerate}
				\item It isn't.
				\item To trace sources of content and information, supporting transparency and possibly legal rights such as copyright.
				\item To increase cost.
				\item To slow inference.
			\end{enumerate}
		\end{block}
		\begin{exampleblock}{Answer: B}\end{exampleblock}
		\begin{block}{Explanation}
			Attribution supports transparency and copyright compliance. Options A, C, and D are incorrect.
		\end{block}
	\end{frame}
	
	\begin{frame}{Q606 --- GenAI and Copyright}
		\begin{block}{Question}
			What is a copyright concern with GenAI?
			\begin{enumerate}
				\item None.
				\item Training on copyrighted material without permission and generating content that may infringe copyright.
				\item GenAI is always copyright-clean.
				\item Copyright does not apply.
			\end{enumerate}
		\end{block}
		\begin{exampleblock}{Answer: B}\end{exampleblock}
		\begin{block}{Explanation}
			Copyright concerns are significant. Options A, C, and D are incorrect.
		\end{block}
	\end{frame}
	
	\begin{frame}{Q607 --- GenAI and Safety}
		\begin{block}{Question}
			Which is a safety concern with GenAI in critical applications?
			\begin{enumerate}
				\item Faster response.
				\item Generating unsafe or harmful content that could cause physical, psychological, or financial harm.
				\item Lower costs.
				\item Higher accuracy.
			\end{enumerate}
		\end{block}
		\begin{exampleblock}{Answer: B}\end{exampleblock}
		\begin{block}{Explanation}
			Harmful content generation is a safety concern. Options A, C, and D are benefits.
		\end{block}
	\end{frame}
	
	\begin{frame}{Q608 --- GenAI and Consumer Protection}
		\begin{block}{Question}
			Which consumer protection issue arises from GenAI?
			\begin{enumerate}
				\item Increasing prices only.
				\item Misleading or inaccurate information provided to consumers, potentially leading to harm.
				\item Faster service.
				\item Lower fees.
			\end{enumerate}
		\end{block}
		\begin{exampleblock}{Answer: B}\end{exampleblock}
		\begin{block}{Explanation}
			Misleading information is a consumer protection issue. Options A, C, and D are unrelated.
		\end{block}
	\end{frame}
	
	\begin{frame}{Q609 --- GenAI and Financial Advice}
		\begin{block}{Question}
			What regulatory concern arises when GenAI is used for financial advice?
			\begin{enumerate}
				\item Lower costs.
				\item Potential non-compliance with advice regulations if unvetted or unsuitable advice is provided.
				\item Faster response.
				\item Higher throughput.
			\end{enumerate}
		\end{block}
		\begin{exampleblock}{Answer: B}\end{exampleblock}
		\begin{block}{Explanation}
			Regulatory concerns exist around advice suitability. Options A, C, and D are benefits, not concerns.
		\end{block}
	\end{frame}
	
	\begin{frame}{Q610 --- GenAI and Fraud}
		\begin{block}{Question}
			How can GenAI facilitate fraud?
			\begin{enumerate}
				\item By making it easier to create convincing phishing content and deepfakes at scale.
				\item By improving accuracy.
				\item By reducing costs.
				\item By speeding up transactions.
			\end{enumerate}
		\end{block}
		\begin{exampleblock}{Answer: A}\end{exampleblock}
		\begin{block}{Explanation}
			GenAI enables scalable, convincing fraud content. Options B, C, and D are unrelated.
		\end{block}
	\end{frame}
	
	% =================================================================
	\section{Emerging Topics}
	% =================================================================
	
	\begin{frame}{Q611 --- Agentic AI Risks}
		\begin{block}{Question}
			What are the key risks of agentic AI?
			\begin{enumerate}
				\item Unintended autonomous actions, difficulty in oversight, and cascading effects from chains of decisions.
				\item Faster decisions.
				\item Lower cost.
				\item Simplicity.
			\end{enumerate}
		\end{block}
		\begin{exampleblock}{Answer: A}\end{exampleblock}
		\begin{block}{Explanation}
			Agentic AI risks relate to autonomy and oversight. Options B, C, and D are benefits or inaccurate.
		\end{block}
	\end{frame}
	
	\begin{frame}{Q612 --- Multi-Agent Systems}
		\begin{block}{Question}
			What risk is specific to multi-agent AI systems?
			\begin{enumerate}
				\item Faster training.
				\item Emergent behaviour arising from interactions between agents, which may be difficult to predict or control.
				\item Lower costs.
				\item Higher accuracy.
			\end{enumerate}
		\end{block}
		\begin{exampleblock}{Answer: B}\end{exampleblock}
		\begin{block}{Explanation}
			Emergent multi-agent behaviour is a specific risk. Options A, C, and D are benefits.
		\end{block}
	\end{frame}
	
	\begin{frame}{Q613 --- AI and Cyber Risk}
		\begin{block}{Question}
			How does AI increase cyber risk?
			\begin{enumerate}
				\item By reducing attack surface.
				\item By enabling more sophisticated, scalable attacks, and by introducing new vulnerabilities (e.g., model extraction, adversarial inputs).
				\item By eliminating threats.
				\item By lowering costs.
			\end{enumerate}
		\end{block}
		\begin{exampleblock}{Answer: B}\end{exampleblock}
		\begin{block}{Explanation}
			AI introduces new attack vectors and scales threats. Options A, C, and D are incorrect.
		\end{block}
	\end{frame}
	
	\begin{frame}{Q614 --- AI Supply Chain Risk}
		\begin{block}{Question}
			What is AI supply chain risk?
			\begin{enumerate}
				\item Risk of low prices.
				\item Risk arising from dependencies on external AI providers, datasets, and libraries, including security, availability, and continuity concerns.
				\item Lower costs.
				\item Higher speed.
			\end{enumerate}
		\end{block}
		\begin{exampleblock}{Answer: B}\end{exampleblock}
		\begin{block}{Explanation}
			Supply chain risk involves external dependencies. Options A, C, and D are benefits or unrelated.
		\end{block}
	\end{frame}
	
	\begin{frame}{Q615 --- AI and Third-Party Risk}
		\begin{block}{Question}
			How should third-party AI risk be managed?
			\begin{enumerate}
				\item Ignore it.
				\item Through due diligence, contractual protections, ongoing monitoring, and contingency planning for provider failures.
				\item Only contracts.
				\item Only internal audits.
			\end{enumerate}
		\end{block}
		\begin{exampleblock}{Answer: B}\end{exampleblock}
		\begin{block}{Explanation}
			Third-party risk requires due diligence, contracts, monitoring, and contingency planning. Options A, C, and D are insufficient.
		\end{block}
	\end{frame}
	
	\begin{frame}{Q616 --- AI and Concentration Risk}
		\begin{block}{Question}
			Why is AI concentration risk a concern?
			\begin{enumerate}
				\item Because it reduces costs.
				\item Because reliance on a small number of AI providers or models creates correlated failure risk across the ecosystem.
				\item Because it increases accuracy.
				\item Because it improves speed.
			\end{enumerate}
		\end{block}
		\begin{exampleblock}{Answer: B}\end{exampleblock}
		\begin{block}{Explanation}
			Concentration amplifies correlated risk. Options A, C, and D are benefits.
		\end{block}
	\end{frame}
	
	\begin{frame}{Q617 --- AI and Reputational Risk}
		\begin{block}{Question}
			Which is a reputational risk from AI?
			\begin{enumerate}
				\item Faster delivery.
				\item Public backlash from biased, unethical, or incompetent AI use, resulting in loss of trust and brand value.
				\item Lower cost.
				\item Higher accuracy.
			\end{enumerate}
		\end{block}
		\begin{exampleblock}{Answer: B}\end{exampleblock}
		\begin{block}{Explanation}
			Reputational risk stems from public perception of AI failures. Options A, C, and D are benefits.
		\end{block}
	\end{frame}
	
	\begin{frame}{Q618 --- AI and Legal Risk}
		\begin{block}{Question}
			Which is a legal risk from AI?
			\begin{enumerate}
				\item Lower cost.
				\item Liability for harm caused by AI decisions, non-compliance with regulations, and IP violations.
				\item Faster deployment.
				\item Higher accuracy.
			\end{enumerate}
		\end{block}
		\begin{exampleblock}{Answer: B}\end{exampleblock}
		\begin{block}{Explanation}
			Legal risks include liability, non-compliance, and IP issues. Options A, C, and D are benefits.
		\end{block}
	\end{frame}
	
	\begin{frame}{Q619 --- AI and Environmental Risk}
		\begin{block}{Question}
			What environmental risk does AI present?
			\begin{enumerate}
				\item None.
				\item Significant compute-related energy consumption and carbon emissions from training and inference of large models.
				\item Lower energy use.
				\item Positive environmental impact only.
			\end{enumerate}
		\end{block}
		\begin{exampleblock}{Answer: B}\end{exampleblock}
		\begin{block}{Explanation}
			AI compute consumption creates environmental risk. Options A, C, and D are incorrect.
		\end{block}
	\end{frame}
	
	\begin{frame}{Q620 --- AI and Regulatory Fragmentation}
		\begin{block}{Question}
			Why is regulatory fragmentation a challenge for AI governance?
			\begin{enumerate}
				\item Because it speeds up compliance.
				\item Because different jurisdictions impose different rules, complicating compliance for globally operating firms.
				\item Because it reduces cost.
				\item Because it simplifies everything.
			\end{enumerate}
		\end{block}
		\begin{exampleblock}{Answer: B}\end{exampleblock}
		\begin{block}{Explanation}
			Fragmented regulation complicates global compliance. Options A, C, and D are incorrect.
		\end{block}
	\end{frame}
	
	% =================================================================
	\section{Integrating AI Governance into the Three Lines}
	% =================================================================
	
	\begin{frame}{Q621 --- First Line AI Responsibilities}
		\begin{block}{Question}
			Which is the first line's responsibility for AI models?
			\begin{enumerate}
				\item Independent validation.
				\item Design, develop, own, and operate AI models; monitor performance and own data quality.
				\item Set policy.
				\item Audit.
			\end{enumerate}
		\end{block}
		\begin{exampleblock}{Answer: B}\end{exampleblock}
		\begin{block}{Explanation}
			First line owns AI models and operations. Options A, C, and D belong to second or third lines.
		\end{block}
	\end{frame}
	
	\begin{frame}{Q622 --- Second Line AI Responsibilities}
		\begin{block}{Question}
			Which is the second line's responsibility for AI models?
			\begin{enumerate}
				\item Build AI models.
				\item Independent validation, policy, and challenge of first-line assumptions and controls.
				\item Use models for decisions.
				\item Market models.
			\end{enumerate}
		\end{block}
		\begin{exampleblock}{Answer: B}\end{exampleblock}
		\begin{block}{Explanation}
			Second line validates, sets policy, and challenges. Options A, C, and D belong to first line.
		\end{block}
	\end{frame}
	
	\begin{frame}{Q623 --- Third Line AI Responsibilities}
		\begin{block}{Question}
			Which is the third line's responsibility for AI models?
			\begin{enumerate}
				\item Build models.
				\item Provide independent assurance on the effectiveness of AI governance and controls.
				\item Set risk appetite.
				\item Manage model data.
			\end{enumerate}
		\end{block}
		\begin{exampleblock}{Answer: B}\end{exampleblock}
		\begin{block}{Explanation}
			Third line provides independent assurance. Options A, C, and D belong to first or second lines.
		\end{block}
	\end{frame}
	
	\begin{frame}{Q624 --- Independence in AI Validation}
		\begin{block}{Question}
			Why must AI validation be independent?
			\begin{enumerate}
				\item To speed up processes.
				\item To avoid conflicts of interest and ensure objective assessment of model quality and risk.
				\item To reduce costs.
				\item To bypass policies.
			\end{enumerate}
		\end{block}
		\begin{exampleblock}{Answer: B}\end{exampleblock}
		\begin{block}{Explanation}
			Independence ensures objective validation. Options A, C, and D are incorrect.
		\end{block}
	\end{frame}
	
	\begin{frame}{Q625 --- Effective Challenge in AI}
		\begin{block}{Question}
			What does effective challenge mean in the AI governance context?
			\begin{enumerate}
				\item Automatic approval.
				\item Critical, objective, and informed review of AI models and their assumptions by qualified parties, independent from developers.
				\item Marketing endorsement.
				\item Rubber-stamping.
			\end{enumerate}
		\end{block}
		\begin{exampleblock}{Answer: B}\end{exampleblock}
		\begin{block}{Explanation}
			Effective challenge is critical review. Options A, C, and D are inconsistent.
		\end{block}
	\end{frame}
	
	\begin{frame}{Q626 --- AI Risk Appetite Cascading}
		\begin{block}{Question}
			How should AI risk appetite cascade through the organisation?
			\begin{enumerate}
				\item It shouldn't.
				\item From board-set appetite to business-unit tolerances and model-level thresholds, all aligned.
				\item Only at business-unit level.
				\item Only at IT level.
			\end{enumerate}
		\end{block}
		\begin{exampleblock}{Answer: B}\end{exampleblock}
		\begin{block}{Explanation}
			Appetite cascades from board to units to models. Options A, C, and D are too narrow.
		\end{block}
	\end{frame}
	
	\begin{frame}{Q627 --- AI Risk Ownership}
		\begin{block}{Question}
			Who owns AI risk in an organisation?
			\begin{enumerate}
				\item The IT department.
				\item Business units that use AI, supported by risk, compliance, and other functions.
				\item Only the CRO.
				\item Only external vendors.
			\end{enumerate}
		\end{block}
		\begin{exampleblock}{Answer: B}\end{exampleblock}
		\begin{block}{Explanation}
			AI risk ownership sits with business users supported by risk functions. Options A, C, and D are incomplete.
		\end{block}
	\end{frame}
	
	\begin{frame}{Q628 --- AI Escalation Paths}
		\begin{block}{Question}
			What is the purpose of AI escalation paths?
			\begin{enumerate}
				\item To slow down decisions.
				\item To ensure that AI incidents, breaches, and issues are reported and addressed at the appropriate level promptly.
				\item To reduce transparency.
				\item To bypass governance.
			\end{enumerate}
		\end{block}
		\begin{exampleblock}{Answer: B}\end{exampleblock}
		\begin{block}{Explanation}
			Escalation paths ensure timely, appropriate responses. Options A, C, and D are contrary.
		\end{block}
	\end{frame}
	
	\begin{frame}{Q629 --- AI Risk Aggregation}
		\begin{block}{Question}
			What does AI risk aggregation mean?
			\begin{enumerate}
				\item Averaging individual risks.
				\item Combining individual AI risks into an enterprise-wide view, accounting for interdependencies and correlations.
				\item Ignoring enterprise risk.
				\item Removing individual risks.
			\end{enumerate}
		\end{block}
		\begin{exampleblock}{Answer: B}\end{exampleblock}
		\begin{block}{Explanation}
			Aggregation considers interdependencies. Options A, C, and D misstate.
		\end{block}
	\end{frame}
	
	\begin{frame}{Q630 --- AI Governance Success Metrics}
		\begin{block}{Question}
			Which is a good metric for AI governance success?
			\begin{enumerate}
				\item Number of meetings.
				\item Reduction in AI incidents, improvements in validation timeliness, and increasing compliance with governance standards.
				\item Number of documents.
				\item Cost of IT.
			\end{enumerate}
		\end{block}
		\begin{exampleblock}{Answer: B}\end{exampleblock}
		\begin{block}{Explanation}
			Governance success metrics track incidents, timeliness, and compliance. Options A, C, and D are not meaningful metrics.
		\end{block}
	\end{frame}
	
	% =================================================================
	\section{Ethics and Regulation in Depth}
	% =================================================================
	
	\begin{frame}{Q631 --- GDPR Data Subject Access Rights}
		\begin{block}{Question}
			Under GDPR, data subjects can request:
			\begin{enumerate}
				\item Only access to data.
				\item Access, rectification, erasure, restriction, portability, and objection.
				\item Only rectification.
				\item Nothing.
			\end{enumerate}
		\end{block}
		\begin{exampleblock}{Answer: B}\end{exampleblock}
		\begin{block}{Explanation}
			GDPR grants multiple rights. Options A, C, and D understate or misstate.
		\end{block}
	\end{frame}
	
	\begin{frame}{Q632 --- GDPR Processor vs Controller}
		\begin{block}{Question}
			What is the difference between a data controller and a data processor under GDPR?
			\begin{enumerate}
				\item They are the same.
				\item Controller determines the purposes and means of processing; processor processes data on behalf of the controller.
				\item Controller only stores data.
				\item Processor only collects data.
			\end{enumerate}
		\end{block}
		\begin{exampleblock}{Answer: B}\end{exampleblock}
		\begin{block}{Explanation}
			Controller decides; processor acts on behalf. Options A, C, and D misstate.
		\end{block}
	\end{frame}
	
	\begin{frame}{Q633 --- GDPR Data Protection Officer}
		\begin{block}{Question}
			What is the DPO's role under GDPR?
			\begin{enumerate}
				\item Marketing.
				\item Independent oversight of data protection compliance, advising on GDPR obligations, and liaising with the supervisory authority.
				\item Sales.
				\item IT support.
			\end{enumerate}
		\end{block}
		\begin{exampleblock}{Answer: B}\end{exampleblock}
		\begin{block}{Explanation}
			The DPO provides independent compliance oversight. Options A, C, and D are unrelated.
		\end{block}
	\end{frame}
	
	\begin{frame}{Q634 --- Data Protection Impact Assessment}
		\begin{block}{Question}
			What is a DPIA?
			\begin{enumerate}
				\item A marketing report.
				\item A Data Protection Impact Assessment required for processing likely to result in high risk to data subjects, documenting risks and mitigations.
				\item A code review.
				\item A financial audit.
			\end{enumerate}
		\end{block}
		\begin{exampleblock}{Answer: B}\end{exampleblock}
		\begin{block}{Explanation}
			DPIA documents high-risk processing and mitigations. Options A, C, and D misstate.
		\end{block}
	\end{frame}
	
	\begin{frame}{Q635 --- AI Act High-Risk Systems}
		\begin{block}{Question}
			Which AI systems are classified as high-risk under the EU AI Act?
			\begin{enumerate}
				\item All chatbots.
				\item Those affecting safety, fundamental rights, or critical infrastructure, such as biometrics, employment, credit scoring, and law enforcement.
				\item Only medical devices.
				\item Only toys.
			\end{enumerate}
		\end{block}
		\begin{exampleblock}{Answer: B}\end{exampleblock}
		\begin{block}{Explanation}
			High-risk systems affect safety, rights, or critical infrastructure. Options A, C, and D are narrower or incorrect.
		\end{block}
	\end{frame}
	
	\begin{frame}{Q636 --- AI Act Transparency Obligations}
		\begin{block}{Question}
			What transparency obligations does the EU AI Act impose?
			\begin{enumerate}
				\item None.
				\item Notifying users when interacting with AI systems, disclosing deepfakes, and providing information about AI-generated content.
				\item Only financial disclosure.
				\item Only marketing disclosure.
			\end{enumerate}
		\end{block}
		\begin{exampleblock}{Answer: B}\end{exampleblock}
		\begin{block}{Explanation}
			The AI Act imposes transparency obligations. Options A, C, and D are incorrect.
		\end{block}
	\end{frame}
	
	\begin{frame}{Q637 --- AI Act General Purpose AI}
		\begin{block}{Question}
			What does the EU AI Act require of general-purpose AI (GPAI) models?
			\begin{enumerate}
				\item Nothing.
				\item Documentation, transparency on training data, copyright compliance, and safety measures, with stricter obligations for models with systemic risk.
				\item Only marketing disclosure.
				\item Only to register.
			\end{enumerate}
		\end{block}
		\begin{exampleblock}{Answer: B}\end{exampleblock}
		\begin{block}{Explanation}
			GPAI models have documentation, transparency, and safety obligations. Options A, C, and D understate or misstate.
		\end{block}
	\end{frame}
	
	\begin{frame}{Q638 --- AI Act Timeline}
		\begin{block}{Question}
			How is the EU AI Act being implemented over time?
			\begin{enumerate}
				\item All at once.
				\item In phases: prohibitions first, then high-risk rules, then remaining provisions, with transitional periods for existing systems.
				\item Only in 2030.
				\item Never.
			\end{enumerate}
		\end{block}
		\begin{exampleblock}{Answer: B}\end{exampleblock}
		\begin{block}{Explanation}
			The AI Act phases in obligations. Options A, C, and D are inaccurate.
		\end{block}
	\end{frame}
	
	\begin{frame}{Q639 --- CCPA Consumer Rights}
		\begin{block}{Question}
			What rights does CCPA grant California consumers?
			\begin{enumerate}
				\item No rights.
				\item Right to know, delete, opt out of sale, correct, and non-discrimination.
				\item Only access.
				\item Only deletion.
			\end{enumerate}
		\end{block}
		\begin{exampleblock}{Answer: B}\end{exampleblock}
		\begin{block}{Explanation}
			CCPA grants multiple rights. Options A, C, and D are incomplete.
		\end{block}
	\end{frame}
	
	\begin{frame}{Q640 --- PIPL Overview}
		\begin{block}{Question}
			What is the purpose of China's Personal Information Protection Law (PIPL)?
			\begin{enumerate}
				\item To ban personal data.
				\item To protect personal data, requiring consent, data minimisation, transparency, and specific rules for cross-border transfer.
				\item Only for government data.
				\item Only for healthcare data.
			\end{enumerate}
		\end{block}
		\begin{exampleblock}{Answer: B}\end{exampleblock}
		\begin{block}{Explanation}
			PIPL protects personal data with consent, minimisation, and transfer rules. Options A, C, and D are inaccurate.
		\end{block}
	\end{frame}
	
	% =================================================================
	\section{Advanced Model Validation}
	% =================================================================
	
	\begin{frame}{Q641 --- Model Validation Evidence}
		\begin{block}{Question}
			What evidence should model validation produce?
			\begin{enumerate}
				\item None.
				\item Documented assessment of conceptual soundness, data appropriateness, implementation, performance, and limitations, with findings and recommendations.
				\item Only accuracy metrics.
				\item Only source code.
			\end{enumerate}
		\end{block}
		\begin{exampleblock}{Answer: B}\end{exampleblock}
		\begin{block}{Explanation}
			Validation produces documented, comprehensive evidence. Options A, C, and D are insufficient.
		\end{block}
	\end{frame}
	
	\begin{frame}{Q642 --- Benchmarking in Validation}
		\begin{block}{Question}
			Why does validation include benchmarking against challenger models?
			\begin{enumerate}
				\item To confuse stakeholders.
				\item To compare performance against alternatives, identifying whether the model is the best choice for its purpose.
				\item To reduce cost.
				\item To speed up training.
			\end{enumerate}
		\end{block}
		\begin{exampleblock}{Answer: B}\end{exampleblock}
		\begin{block}{Explanation}
			Benchmarking compares models for suitability. Options A, C, and D are incorrect.
		\end{block}
	\end{frame}
	
	\begin{frame}{Q643 --- Sensitivity Analysis}
		\begin{block}{Question}
			What does sensitivity analysis in model validation examine?
			\begin{enumerate}
				\item Only cost.
				\item How model outputs respond to changes in input assumptions, parameters, or data, evaluating robustness.
				\item Only speed.
				\item Only storage.
			\end{enumerate}
		\end{block}
		\begin{exampleblock}{Answer: B}\end{exampleblock}
		\begin{block}{Explanation}
			Sensitivity analysis tests robustness to changes. Options A, C, and D are unrelated.
		\end{block}
	\end{frame}
	
	\begin{frame}{Q644 --- Stress Testing AI Models}
		\begin{block}{Question}
			What is the purpose of stress testing AI models?
			\begin{enumerate}
				\item To make them faster.
				\item To evaluate their behaviour under extreme, adverse, or unlikely scenarios.
				\item To reduce cost.
				\item To simplify them.
			\end{enumerate}
		\end{block}
		\begin{exampleblock}{Answer: B}\end{exampleblock}
		\begin{block}{Explanation}
			Stress tests evaluate behaviour under extremes. Options A, C, and D are unrelated.
		\end{block}
	\end{frame}
	
	\begin{frame}{Q645 --- Backtesting AI Models}
		\begin{block}{Question}
			What does backtesting involve for AI models?
			\begin{enumerate}
				\item Forward-looking simulation only.
				\item Testing the model against historical data to evaluate how it would have performed.
				\item Marketing analysis.
				\item Cost analysis.
			\end{enumerate}
		\end{block}
		\begin{exampleblock}{Answer: B}\end{exampleblock}
		\begin{block}{Explanation}
			Backtesting evaluates historical performance. Options A, C, and D are unrelated.
		\end{block}
	\end{frame}
	
	\begin{frame}{Q646 --- Outcome Analysis}
		\begin{block}{Question}
			Outcome analysis in model validation is used to:
			\begin{enumerate}
				\item Reduce costs.
				\item Compare model predictions against realised outcomes to detect bias, degradation, or incorrect assumptions.
				\item Speed up training.
				\item Increase accuracy.
			\end{enumerate}
		\end{block}
		\begin{exampleblock}{Answer: B}\end{exampleblock}
		\begin{block}{Explanation}
			Outcome analysis compares predictions to reality. Options A, C, and D are unrelated.
		\end{block}
	\end{frame}
	
	\begin{frame}{Q647 --- Use Test Evidence}
		\begin{block}{Question}
			What evidence should the use test provide?
			\begin{enumerate}
				\item None.
				\item Qualitative evidence that the model is used appropriately, understood by users, and aligned with business decisions.
				\item Only accuracy metrics.
				\item Only cost data.
			\end{enumerate}
		\end{block}
		\begin{exampleblock}{Answer: B}\end{exampleblock}
		\begin{block}{Explanation}
			The use test produces qualitative evidence. Options A, C, and D are insufficient.
		\end{block}
	\end{frame}
	
	\begin{frame}{Q648 --- Model Use Restrictions}
		\begin{block}{Question}
			How should model use restrictions be enforced?
			\begin{enumerate}
				\item Through documentation only.
				\item Through documentation, access controls, system limits, and monitoring.
				\item By ignoring them.
				\item By removing them.
			\end{enumerate}
		\end{block}
		\begin{exampleblock}{Answer: B}\end{exampleblock}
		\begin{block}{Explanation}
			Use restrictions require documentation, controls, limits, and monitoring. Options A, C, and D are insufficient or contrary.
		\end{block}
	\end{frame}
	
	\begin{frame}{Q649 --- Model Validation Documentation}
		\begin{block}{Question}
			What documentation should accompany model validation?
			\begin{enumerate}
				\item Nothing.
				\item A validation report with methodology, findings, limitations, and recommendations, plus supporting analyses.
				\item Only the model code.
				\item Only marketing material.
			\end{enumerate}
		\end{block}
		\begin{exampleblock}{Answer: B}\end{exampleblock}
		\begin{block}{Explanation}
			Validation documentation is comprehensive. Options A, C, and D are insufficient or inappropriate.
		\end{block}
	\end{frame}
	
	\begin{frame}{Q650 --- Responsible AI Closing Principle}
		\begin{block}{Question}
			Which best summarises the guiding principle for responsible AI governance?
			\begin{enumerate}
				\item Maximise complexity.
				\item Balance performance with governance, validation, monitoring, fairness, explainability, and ethical alignment, consistent with regulatory and business expectations.
				\item Deploy without validation.
				\item Ignore risks.
			\end{enumerate}
		\end{block}
		\begin{exampleblock}{Answer: B}\end{exampleblock}
		\begin{block}{Explanation}
			Responsible AI balances performance and governance across all dimensions. Options A, C, and D contradict the principle.
		\end{block}
	\end{frame}
	
\end{document}